% Normative Framework: Process vs. Outcome Legitimacy
% This is the conceptual core for Section 6.X.4 - addressing whether
% geography-induced bias is ethically different from intentional gerrymandering

\subsubsection{The Normative Question: Process Legitimacy vs. Outcome Proportionality}
\label{sec:normative_framework}

The empirical findings raise a fundamental normative question: If algorithmic redistricting produces systematically biased partisan outcomes due to geographic sorting, is this ethically different from intentional gerrymandering?

\paragraph{Three Positions on Geography-Induced Bias}

\textbf{Position 1: Geography-Induced Bias is Legitimate (Geographic Realism).}
This position holds that partisan bias arising from residential patterns can be normatively acceptable. Any districting system respecting geographic constraints will reflect some of those patterns \citep{rodden2019why}. Requiring proportional outcomes may demand tradeoffs with compactness or community boundaries; whether those tradeoffs count as manipulation is itself contested.

\textit{Constitutional support}: \textit{Rucho v. Common Cause} explicitly rejected proportional representation requirements. Chief Justice Roberts: ``Deciding among different visions of fairness poses basic questions that are political, not legal'' \citep{rucho2019}.

\textit{Critique}: This may be overly accepting of consistent minority rule. If geographic sorting systematically advantages one party across many states, cumulative national effects could be anti-democratic even if individual state processes are neutral.

\textbf{Position 2: All Bias is Problematic (Outcome-Based Fairness).}
This position holds that partisan bias---regardless of source---violates political equality norms. Proponents argue that seat shares should track vote shares more closely \citep{mcgann2016gerrymandering}. Redistricting should actively compensate for geographic disadvantage through efficiency-gap targets, proportional representation mechanisms, or symmetry requirements.

\textit{Constitutional challenge}: This conflicts with \textit{Rucho}, which rejected outcome-based standards as non-justiciable political questions. Courts cannot determine ``what is fair'' because ``fairness'' is inherently contested.

\textit{Critique}: Requires contested normative judgments about fair representation. Should we target proportionality? Competitiveness? Symmetry? Each rests on different assumptions and produces different results.

\textbf{Position 3: Process Legitimacy (Our Position).}
Geography-induced bias is legitimate \textit{if and only if} the redistricting process was neutral and did not exploit geographic patterns for partisan advantage. Process fairness takes priority over outcome proportionality.

\textit{Core argument}: One account grounds democratic legitimacy partly in procedural fairness \citep{rawls1971theory}. On this view, transparency, equal treatment, public criteria, and reproducibility contribute to legitimacy but do not by themselves settle legality or substantive fairness.

\textit{Key distinction}: The normative difference lies in bias \textit{source}:
\begin{itemize}
\item \textbf{Intentional gerrymandering}: Geography is a \textit{tool} for manipulation
\item \textbf{Algorithmic redistricting}: Geography is a \textit{constraint}, not a tool
\end{itemize}

\textit{Legal alignment}: This threads the needle in \textit{Rucho}. Roberts rejected outcome-based standards but invited ``neutral criteria such as compactness.'' Process-based standards are objectively verifiable---courts can determine whether partisan data was used (yes/no) without judging whether outcomes are ``fair enough'' (contested).

\paragraph{Justification for Process Legitimacy}

We adopt Position 3 on four grounds:

\textbf{First, legal realism.} \textit{Rucho} forecloses outcome-based standards in federal courts. Any viable reform post-\textit{Rucho} must focus on process. State constitutional litigation has validated this approach (Pennsylvania struck maps for violating compactness, not for biased outcomes).

\textbf{Second, philosophical coherence.} Process legitimacy draws on accounts of procedural justice \citep{rawls1971theory}. In redistricting, however, a verifiable process is evidence about execution rather than a guarantee that every output is substantively fair.

\textbf{Third, practical implementability.} Process-based standards are judicially manageable: objective verification (was partisan data used?), clear ex ante guidance (use compactness, ignore partisanship), and uniform application across states.

\textbf{Fourth, normative acceptability of geographic constraints.} Geography is not infinitely malleable. Requiring proportional representation in Massachusetts (extreme urban concentration) would necessitate radical non-compactness or combining incompatible communities. Accepting geographic reality respects federalism and local communities.

\paragraph{Implications for Algorithmic Redistricting}

Geography-induced bias is acceptable when: (1) the algorithm optimizes neutral criteria without partisan data access, (2) partisan outcomes are byproducts of geographic optimization, (3) the process is transparent and reproducible, and (4) outcomes match geographic expectations. Our approach satisfies all four---the algorithm has zero partisan degrees of freedom.

The key question is not ``Are outcomes proportional?'' but ``Could the process have produced different outcomes through strategic manipulation?'' For algorithmic redistricting, the answer is demonstrably no.

Traditional redistricting often fails these tests. Even ``bipartisan commissions'' engage in strategic bargaining using geographic knowledge. The process is neither transparent nor reproducible---no one can verify whether decisions were neutral or strategically motivated.

\paragraph{Addressing Cumulative National Effects}

A remaining challenge: if geographic sorting systematically advantages Republicans across many states, aggregate effects in the House could produce persistent asymmetric bias.

This concern has force but points toward different remedies: (1) Geography's effect varies by decade and is historically contingent, not permanent \citep{chen2017impact}. (2) If cumulative bias is the concern, remedies should operate at the \textit{national} level (expanding the House, multi-member districts) rather than distorting \textit{state-level} processes. (3) Parties suffering geographic disadvantage have political recourse: change platforms to appeal to different voters or pursue constitutional amendments.

The strongest version of process legitimacy accepts that some geographic configurations advantage one party. This is a feature of single-member district systems in geographically sorted societies. The appropriate response is ensuring redistricting processes are scrupulously neutral, then letting political competition adapt to geographic realities.

\paragraph{Conclusion: Geography as Constraint, Not Tool}

We distinguish between:
\begin{itemize}
\item \textbf{Geography as constraint}: Residential patterns that redistricting must respect (legitimate)
\item \textbf{Geography as tool}: Strategic use of geographic knowledge for desired outcomes (illegitimate)
\end{itemize}

Algorithmic redistricting treats geography purely as constraint---the algorithm has no knowledge of partisan patterns and cannot exploit them. This is precisely what \textit{Rucho} invited: redistricting based on ``neutral criteria such as compactness'' rather than partisan outcome targets.

The question is not whether algorithmic maps achieve perfect proportionality (they do not, because geography does not permit it), but whether they eliminate partisan manipulation (they do, because the algorithm has zero partisan degrees of freedom). Process legitimacy demands the latter, not the former.
